\documentclass[letterpaper,11pt,preprint]{aastex}
\usepackage{graphics,graphicx}
\usepackage{amsmath}
\usepackage{natbib}
\usepackage{color}
\citestyle{aas}

\setlength{\textwidth}{6.5in} \setlength{\textheight}{9in}
\setlength{\topmargin}{-0.0625in} \setlength{\oddsidemargin}{0in}
\setlength{\evensidemargin}{0in} \setlength{\headheight}{0in}
\setlength{\headsep}{0in} \setlength{\hoffset}{0in}
\setlength{\voffset}{0in}

\makeatletter
\renewcommand{\section}{\@startsection%                                                                                                                                                                    
{section}{1}{0mm}{-\baselineskip}%                                                                                                                                                                         
{0.5\baselineskip}{\normalfont\Large\bfseries}}%                                                                                                                                                           
\makeatother

%%%%%%%%%%%%%%%%%%%%%%%%%%%%%                                                                                                                                                                              
%%%%% Start of document %%%%%                                                                                                                                                                              
%%%%%%%%%%%%%%%%%%%%%%%%%%%%%                                                                                                                                                                              

\begin{document}
\pagestyle{plain}

\section{Pressure and the Ideal Gas Law}

If you were to zoom in on a little piece of air, you would see a bunch
of molecules, mostly nitrogen and oxygen, bouncing around randomly.
The molecules would be spaced relatively far apart, occupying only about
0.1\% of the volume of the air, and 99.9\% empty space.  The molecules
regularly smack into each other and bounce off, with each individual
molecule getting hit billions of times per second.  When we zoom back
out, somehow this messy, chaotic system turns into something quite
orderly that can transport music from an intrument to our ears.  We're
going to try to understand at least a little bit about how this
happens, and, eventually (but not in this note), how this messy
microscopic picture actually can affect what we hear.

The first concept we want to work out is \textit{pressure}.  Pressure
is defined to be the force per unit area a medium exerts on a
surface.  All sorts of things can cause pressures in physics,
including light, magnetic fields, even quantum mechanics.  Air and
water pressures are the ones we most commonly experience directly, and
in this class we'll just look at air pressure.

\textit{Example: Pressure on a Jar of Pickles}\newline
To get a sense of how pressure behaves, let's try to open a jar of
pickles.  The pressure at sea level is about 100,000 Pascals, or 100
kPa, where a Pascal is a Newton per square metre.  This isn't really a
very intuitive number, so let's break it down.  Nobody I know actually
thinks in Newtons; have you ever picked up a suitcase and said
``this weighs about 200 Newtons''?  Me neither.  We can convert
between mass and weight using Newton's second law:
\begin{equation}
  F=ma
\end{equation}
Here, $F$ is the force, $m$ is the mass, and $a$ is the acceleration,
the rate of change of velocity.  Let's also be careful about units -
$m$ is in $kg$, a velocity is in $m/s$, so a \textit{change} in
velocity per unit time is $(m/s)/s$ or $m/s^2$.  In an equation, the
units on both sides have to match, so the units of force must be the
units of mass times the units of acceleration, or $kg-m/s^2$.  This is
a Newton - the force required to make a mass of 1 $kg$ speed up by 1
$m/s$ each second.  The acceleration from Earth's gravity $g$ at the
surface of the Earth is about $9.8 m/s$, which we'll round to $10 m/s$.
If we want the gravitational force, we can swap in $g$ for $a$ in
Newton's second law, giving $F_g=mg$.  If we have a weight of 1 $kg$
then, it takes a force of about 10$N$ to counter the force from
gravity.

If we go back to atmospheric pressure, then 100,000$N$ is the force
required to hold up $100,000/10=10,000 kg$ or 10 tons.  That's a
\textit{huge} number, so why don't all the windows in your house blow
out?  Well, the pressure on the inside of your house is the same as
the outside, so the forces from the pressures cancel out.  If the
pressure inside your house was lower, though, your windows would feel
a big force and would most likely fail spectacularly.

The pressure inside a pickle jar is indeed much lower than
atmospheric.  If you've ever canned anything, you know that you heat
up whatever you're canning, put it in a jar, seal it up, and let it
close.  When gasses cool, their pressure drops, so the pressure inside
a jar is around 20-25\% lower than atmospheric.  Let's say the
diameter of the jar lid is 10$cm$ or 0.1$m$, so the area is $\pi r^2
\sim 3 (0.1/2)^2 \sim 0.008m^2$.  The force is from the amosphere is
then $0.008 \times 10^5 = 800N$ downwards, or about an 80 $kg$ weight on top of
the lid.  If the pressure inside is 75\% of atmospheric, then there's
an upward force of $600N$, so the total force on the lid is 200$N$,
equivalent to 20$kg$.  When you unscrew the lid, you have to counter
that force, which is why it's hard to open a jar the first time.  Once
you're open, though, the pressures are the same, and the next time you
open is much easier!

\section{Pressure from a Gas}
Now that we (hopefully!) have at least a little bit of sense for what
pressure feels like, let's look at where it comes from.  If we zoom in
on our pickle jar lid, we see a bunch of air molecules zipping around,
bouncing off the lid.  Every time a molecule bounces off the lid, its
velocity changes\footnote{``But wait!'' you say.  ``If I bounce off
the lid, my speed is the same before and after I bounce!''  Yes,
that's true, but velocity is speed \textit{and direction}.  If I move
left, bounce, then move right at the same speed, my velocity changed
because my direction changed.}, and a change in velocity means an acceleration.
Newton's second law tells us that means the lid had to exert a force
on the molecule.  For a single molecule, the force is tiny, but there
are a \textit{huge} number of collisions each second, and their
combined effects are not tiny.  Let's try to work out what this looks
like now.

We'll start by considering a long, thin tube (like a very, very long
straw).  The tube is closed at the ends, and the area of each end
we'll call $A$.  We've got a bunch of molecules bouncing around inside
our tube, and we'll try to work out the force on the end.  The first
thing to notice is that we don't care about up/down or front/back
motions if we're only looking at the end.  Let's say we're looking at
the left end - then any molecule near it moving to the left will
bounce, and it will start moving to the right.  Picture a bounce pass
in basketball - the horizontal speed doesn't change, and the vertical
speed flips from down to up.

Now that we know what we're looking for, let's try to work out the
force.  We're going to estimate a total mass bouncing off the end and
the change in velocity, then plug those into Newton's second law to
get a force.  Let's call the ``typical'' speed in the $x-$direction $v_x$,
the typical mass of each molecule $m$, and the number of molecules per
unit volume $n$.  We're going to make the simplifying assumption that
half the molecules are moving to the left with $v_x$ and half to the
right.  What happens at the left end cap?  If we wait for a
length of time $t$, then half of our molecules will move a distance
$v_xt$ to the left, and half will move the same distance to the
right.  If we ignore collisions and assume each air molecule is just
doing its own thing, then during our time interval $t$, every air
molecule moving to the left within $v_xt$ of the end cap will bounce
off the cap.  Because the area of our  tube is $A$, then every
left-moving molecule with a volume of $v_xtA$ is going to bounce off.
The number of molecules is going to be the number of things moving
left per volume which is $n/2$ times the volume, so we have $v_x
tAn/2$ collisions per second.  The total mass that bounces off our end
cap is the number of objects that hit the end times the mass per object, or
$mv_xtAn/2$.  The \textit{change} in velocity is going to be $2v_x$,
which happens over a time $t$, so the acceleration is $2v_x/t$.  The
force is now $mv_xtAn/2 \times 2v_x/t = mv_x^2 An$.  To get the force
per unit area, we take the force and divide by the area $A$, leaving
us with the simple expression:
\begin{equation}
  P=mv_x^2 n
  \label{eqn:ideal_gas_kinetic}
\end{equation}
Now, you might remember that $E=1/2 mv^2$ is the kinetic energy of an
object.  We can rewrite the previous equation to be $P=2(1/2 mv_x^2) n
= 2E_x n$ where $E_x$ is the kinetic energy from the velocity in the
$x-$direction.  What is that energy, though?  Well, you may remember
being told that ``temperature is a measure of the average energy''.
Energy is in Joules, temperature (if you're a physicist) is in Kelvin,
and in our case we can go back and forth between them using
\begin{equation}
  E_x=\frac{1}{2}k_BT
  \label{eqn:Boltzmann}
\end{equation}
where $T$ is temperature and $k_B$ is one of the fundamental constants
of nature, Boltzmann's constant.  The factor of $1/2$ comes from
statistical mechanics, and its derivation is well beyond the scope of
this class, but I hope you don't mind taking the result on faith.

Once we have Equation~\ref{eqn:ideal_gas_kinetic}, we plug in
Equation~\ref{eqn:Boltzmann}, and we have the physicists' version of
the ideal gas law:
\begin{equation}
  P=nkT
  \label{eqn:ideal_gas_one}
\end{equation}
We can go halfway to a chemist's version by noting if we have $N$
things in a volume $V$, that $n=N/V$.  Plug that into
Equation~\ref{eqn:ideal_gas_one} and we have
\begin{eqnarray}
  P=N/V kT\\
  PV=NkT
\end{eqnarray}
In this case, $N$ is the number of molecules we have in our volume
$V$, but Chemists like to work with moles of things.  I need
Avogadro's number $N_A=6.02e23$ of things to get one mole.  I can
stick in two cancelling copies of $N_A$ into the previous equation to
get:
\begin{eqnarray}
  PV=N (N_A/N_A)kT\\
  PV=(N/N_A) (kN_A)T\\
  PV=N_m rT
\end{eqnarray}
where $r\equiv kN_A$ is known as the ideal gas constant, and $N_m$ is
the number of moles of stuff we have.  Boltzmann's constant is
$1.38e-23 J/K$, and so $r=1.38e-23 J/K \times 6.02e23/mole = 8.31
J/K/mole$ -- a number you may remember from chemistry class (if you
took chemistry).

Before doing the recap, I'd like to make an aside, in case any of you
are still worried about this derivation.  Feel free to skip this,
however.  Pretending every molecule is
moving with the same speed either in the $+x$ or $-x$ direction seems
like a pretty big cheat.  And it is!  In reality, the molecules will
all have random velocities, so to do this right, we'd need to average
over those velocities.  If you look at
Equation~\ref{eqn:ideal_gas_kinetic}, the right thing to do is to
split up all our molecules into groups that are moving with the same
velocity.  We then sum over all those little velocity slices.  Forces
add, so pressures have to add, which leaves us with
$$P=\sum_{slice} m v_{x,slice}^2n_{slice}$$
Because the mass is the same for everybody, we can pull it out of the
sum.  We can also define $f_{slice}=n_{slice}/n$, that is $f_{slice}$
is the \textit{fraction} of molecules with velocity $v_{slice}$.  That
turns our sum into:
\begin{equation}
  P=mn\sum_{slice}v_{x,slice}^2f_{slice}
\end{equation}
That sum is exactly what we mean by the average of something, so in
this case, what we should have used was the average value of $v_x^2$,
which we denote with angle brackets: $\left < v_x^2\right >$. That
leaves us with
\begin{equation}
  P=mn\left <v_x^2 \right > = 2n \left <m v_x^2/2 \right >
\end{equation}
That expression is just the average kinetic energy from the
$x-$velocity, which we know is $\frac{1}{2}kT$ (well, you would know
if you had taken a course in statistical mechanics, but if you
haven't, just take my word for it).  When we made our seemingly
horribly crude approximation that all molecules move with the same
speed, we would generally expect that we'd get a slightly wrong (but
close!) answer.  In this case, though, once we saw what our answer
was, we can actually see that we got exactly the right answer because
the average we would need to do is exactly the average statistical
mechanics already did for us.

We've now derived from first principles, only using one extra bit of
information from statistical mechanics, that if I put a gas inside a
container, the container will feel a pressure from the gas.  The more
stuff I put in the container, the higher the pressure will be.  The
hotter the gas, again, the higher the pressure will be.  You'll notice
that the mass of the things dropped out, so it doesn't matter
\textit{what} kind of gas I put in, just how many molecules.

One final note - you might say ``Sure, this gives us the
\textit{average} force, but in reality it's a bunch of molecules
hitting a wall.  Does that matter?''  It might!  But probably not.
Let's take the pickle jar again.  The typical $x-$velocity is
$\frac{1}{2}mv_x^2=\frac{1}{2}kT$ which we can solve to get
\begin{equation}
  v_x=\sqrt{\frac{kT}{m}}
\end{equation}
The mass of a mole of protons/neutrons is around 1
gram\footnote{Originally, chemists wanted a mole of oxygen to be 16
grams, while physicists wanted a mole of hydrogen to be one gram.
Much like Ottawa as the capitol of Canada was a compromise between
Toronto and Montreal, physicists and chemists comprimised and decided
that a mole would be 12 grams of carbon-12.}.  The
average atomic weight of air is around 30 (nitrogen is 28, oxygen is
32), so the mass of an air molecule is $30/N_A\ g=30/N_A/1000kg \sim
5e-26 kg$.  We have to measure room temperature in Kelvin, so add 273
to temperature in Celsius to get $T\sim 293K$.  That gives a typical
speed in the $x-$direction of
\begin{equation}
  |v_x|\sim \sqrt{\frac{1.38e-23\times 293}{5e-26} }=280 m/s
\end{equation}
If we count the collisions in one second, our pickle jar lid then
picks up half the molecules in a volume of $280 m/s *0.008 m^2 *1 s
\sim 2.2 m^3 = 2200l$.  The density of air is about one gram per
litre, so 2200 grams of air or 2.2 $kg$ bounce off our pickle lid.  We
know that the mass of an air molecule is $5e-26 kg$, so we have
$2.2/5e-26=4.4e25$ air molecules bouncing off our pickle jar lid every
second.  Because the air molecules are moving randomly, some seconds
we'll get a few more bouncing off of us, some seconds a few less.  I
won't derive it here, but the typical scatter in that goes like the
square root of the number of things, so we'd expect
$\sqrt{4.4e25}=6.6e12$.  While that may seem like a large number (6
trillion more air molecules hit me this second than hit me last
second!), because the total number of collisions is so huge, you
almost certainly would never notice it.  What we can measure is the
force on the lid, and the \textit{fraction} change in that is
$6.6e12/4.4e25=1.5e-13$.  That is, the randomness in the force on our
jar lid is less than a part per trillion.  That tiny change in force
is so small you'd never notice it.  However, if you zoom way in and
look on very short times, the randomness \textit{isn't} tiny, and you
can see very small particles jitter around in what is
called Brownian Motion.  A young scientist worked out much of the math
behind this in 1905 in order to predict the (random) motions of small
objects buffetted by collisions with individual gas or water
molecules.  That scientist's name?  Albert Einstein.



\end{document}
